\documentclass[10pt,a4paper]{amsart}
\usepackage[margin=19mm]{geometry}
\usepackage{amsmath,amssymb,mathtools}
\usepackage[scr=rsfs,cal=euler]{mathalfa}
\usepackage{xfrac}
\usepackage[unicode,hidelinks,bookmarks=false]{hyperref}
\newcommand{\ie}{i.e.\ }
\newcommand{\cf}{cf.\ }
\newcommand{\resp}{respectively\ }
\newcommand{\Subsection}{\subsection}
\providecommand{\nobreakdash}{\nobreak\hbox{-}}
\setlength{\emergencystretch}{2em}
\multlinegap0pt
\usepackage{amssymb}
\usepackage{tikz}
\usepackage{float}
\newtheorem{thm}{Theorem}
\newtheorem{cor}{Corollary}
\title{Regular sets of circulant quartic graphs}
\author{A. Abdollahi, J.Bagherian, F. Jafari, M. Khatami, Z. Shokoohi, R. Sobhani}
\date{Source: arXiv:2609.09414v1 (CC BY 4.0)}
\begin{document}
\maketitle

\noindent\emph{Source and licence.} Regular sets of circulant quartic graphs. Version arXiv:2609.09414v1 (CC BY 4.0), distributed by arXiv under the Creative Commons Attribution 4.0 International licence, http://creativecommons.org/licenses/by/4.0/, as recorded in the arXiv licence metadata for this exact version and shown on the abstract page https://arxiv.org/abs/2609.09414v1. Full licence notice: \texttt{licenses/arxiv-2609.09414-CC-BY-4.0.txt}.\par\smallskip

\noindent\textit{Editorial scope.} Counted unit: the article's cor cor:coset-size2 (source0.tex lines 419--437). The article's complete original proof of it is delivered with it (source0.tex lines 438--520, one \texttt{proof} environment of 83 lines). No mathematics is written, repaired or re-derived: the statement and the proof are the article's own lines, in the article's own notation, and no retained line is substituted or re-flowed.  Retained uncounted as the source-local context the proof needs: lines 294--298.  Nothing was omitted from any retained span, so the package carries no omission record.  No retained line contains a citation, so the wrapper carries no bibliography and no bibliography entry.  No retained line contains a figure environment, an image-inclusion command or drawing code, so no image is delivered and the package declares no asset dependency.  Editorial formatting only, in addition: an AMS \texttt{amsart} preamble that restates the article's own declarations and macros used by the retained lines, namely the environments \texttt{thm}, \texttt{cor} on separate counters, together with the standard packages those lines need.  The retained environments print in this document as Theorem 1, Corollary 1 in order of appearance, whereas the article prints the same items with the numbering of its own full document.  The complete native source of the article as distributed in the arXiv e-print for this exact version is bundled unchanged as \texttt{sources/209812/source0.tex}.
	\begin{thm}\label{khi3}
		Let $n$ be an integer and $p$ be an prime number such that  $p \mid n$. Also let $C$ be an $(a,b)$-regular set in $\mathrm{Cay}(\mathbb{Z}_n, S)$, $s_0 := \lvert \langle x^p \rangle \cap S \rvert + b - a$ and
		$s_d := \lvert x^d \langle x^p \rangle \cap S \rvert$ for each integer $d \in [1,\, p-1]$. If $\sum_{d=0}^{p-1}s_d\neq 0$, and there exist integers, $0 \leq u,v \leq p-1$ such that $s_u \neq s_v$, then $p \mid |C|$ and
		$\lvert x^h \langle x^p \rangle \cap C \rvert = \frac{|C|}{p}$ for all \(0 \le h \le p-1\).
	\end{thm}

\begin{cor}\label{cor:coset-size2}
	Let $n$, $k$, $\ell$, and $p$ be positive integers, where $p$ is a prime divisor of $n$. Let $\Gamma=\operatorname{Cay}(\mathbb{Z}_n,\{x^k,x^{-k},x^\ell,x^{-\ell}\})$ be a connected circulant quartic graph, and $C$ be an $(a,b)$-regular set in $\Gamma$ such that $p\nmid k$ and $p\nmid \ell$. If one of the following conditions holds:
	\begin{enumerate}
		\item[(i)] $p=2$ and $b-a\notin\{-4,4\}$;
		\item[(ii)] $p=3$ and $b-a\notin\{-4,2\}$;
		\item[(iii)] $p=5$, $b-a\neq -4$, and
		$k\equiv \pm \, r \pmod{5}$ and $\ell\equiv \pm \, r \pmod{5}$,
		where $r\in\{1,2\}$;
		\item[(iv)] $p=5$, $b-a\notin\{-4,1\}$, and $k\equiv \pm r_1 \pmod{5}$ and $\ell\equiv \pm r_2 \pmod{5}$,
		where $r_1,r_2\in\{1,2\}$ are distinct;
		\item[(v)] $p>5$ and $b-a\neq -4$,
	\end{enumerate}
	then
	\[
	\left|x^h\langle x^p\rangle\cap C\right|
	=\frac{bn}{p(4+b-a)}
	\]
	for every $0\le h\le p-1$;
\end{cor}

\begin{proof}
	By substituting
	$
	S=\{x^k,x^{-k},x^\ell,x^{-\ell}\}
	$
	into Theorem~\ref{khi3}, we consider the following cases for
	$
	(\overline{S}+b-a)^{\varphi_p}.
	$ We then verify the hypotheses of Theorem~\ref{khi3}; namely,
		\begin{equation}\label{eq:susv2}
		\exists\,u,v\in\{0,\ldots,p-1\}\ \text : \ s_u \neq s_v, \end{equation}
	and $\sum_{d=0}^{p-1}s_d\neq0.$	
	
	\indent\textbf{(i)} Suppose that \(p=2\). Since $2\nmid k$ and $2\nmid \ell$, we have $(x^k+x^{-k})^{\varphi_2}=2X$,  $(x^\ell+x^{-\ell})^{\varphi_2}=2X$, where $X=x^\frac{n}{2}$. Hence,
	\[
	(\overline{S}+b-a)^{\varphi_2}=(b-a)+4X.
	\]
We have
$s_0=b-a$ and $s_1=4.$
By the assumption in part~(i), \(b-a\neq 4\), and hence
$s_0\neq s_1.$
Moreover, \(b-a\neq -4\), and therefore $s_0+s_1\neq 0.$ So,
$
\left|x^h\langle x^2\rangle\cap C\right|
=
\frac{bn}{2(4+b-a)}$ for $h=0,1.$

\indent\textbf{(ii)} Suppose that \(p=3\), and $3\nmid k$ and $3\nmid \ell$. We have $(x^k+x^{-k})^{\varphi_3}=X+X^{-1}$, $(x^\ell+x^{-\ell})^{\varphi_3}=X+X^{-1}$, where $X=x^\frac{n}{3}$. Hence,
\[
(\overline{S}+b-a)^{\varphi_3}=(b-a)+2X+2X^{-1}.
\]	
We have
$s_0=b-a$ and $s_1=s_2=2.$ Since \(b-a\neq 2\), we have \(s_0=b-a\neq2=s_1=s_2\). Thus, there exist \(u,v\in\{0,1,2\}\) such that \(s_u\neq s_v\). Moreover, since \(b-a\neq -4\), it follows that \(\sum_{i=0}^{2}s_i\neq 0\).
So, $
\left|x^h\langle x^3\rangle\cap C\right|
=
\frac{bn}{3(4+b-a)}$ for $h=0,1,2.$   	

\indent\textbf{(iii)} Since $p=5$, and $k\stackrel{5}{\equiv} \pm\, r$ and   $\ell\stackrel{5}{\equiv} \pm\, r$ for some $r\in\{1,2\}$, then
\[
(\overline{S}+b-a)^{\varphi_5}=(b-a)+2X^{r}+2X^{-r}
\]
 Where $X=x^\frac{n}{5}$. Since $s_0=b-a$, exactly two of the numbers $s_1,\ldots,s_{4}$ are equal to $2$, and the remaining numbers are equal to $0$. Hence, there exist distinct indices $u,v\in\{1,\ldots,4\}$ such that $s_u\neq s_v$. Moreover, since \(b-a\neq -4\), it follows that \(\sum_{i=0}^{4}s_i\neq 0\).
So, $
\left|x^h\langle x^5\rangle\cap C\right|
=
\frac{bn}{5(4+b-a)}$ for $0\le h\le 4$.

\indent\textbf{(iv)} Suppose that $p=5$. Let $X=x^\frac{n}{5}$. Since $k\stackrel{5}{\equiv} \pm\, r_1$ and $\ell\stackrel{5}{\equiv} \pm\, r_2$, where $r_1,r_2\in\{1,2\}$ are distinct, we have $(x^k+x^{-k})^{\varphi_5}=X^{r_1}+X^{-r_1}$ and $(x^\ell+x^{-\ell})^{\varphi_5}=X^{r_2}+X^{-r_2}$. Hence,
\[
(\overline{S}+b-a)^{\varphi_5}
=(b-a)+X^{r_1}+X^{-r_1}+X^{r_2}+X^{-r_2}
\]
We have
$s_0=b-a$ and $s_1=s_2=s_3=s_4=1.$ Since \(b-a\neq 1\), we have \(s_0=b-a\neq1=s_1=s_2=s_3=s_4\). Thus, there exist \(u,v\in\{0,1,2,3,4\}\) such that \(s_u\neq s_v\).  Moreover, since \(b-a\neq -4\), it follows that \(\sum_{i=0}^{4}s_i\neq 0\).
So, $
\left|x^h\langle x^5\rangle\cap C\right|
=
\frac{bn}{5(4+b-a)}$ for $0\le h\le 4$.

\indent\textbf{(v)} Since $p\nmid k$ and $p\nmid \ell$, there exist $r_1,r_2\in\left\{1,\ldots,\frac{p-1}{2}\right\}$ such that $k\stackrel{p}{\equiv}\pm r_1$ and $\ell\stackrel{p}{\equiv}\pm r_2$. Let $X=x^{n/p}$. If $r_1=r_2$, then
\begin{equation}\label{eq:sbar1}
(\overline{S}+b-a)^{\varphi_p}=(b-a)+2X^{r_1}+2X^{-r_1}	
\end{equation}
Otherwise, if $r_1\neq r_2$, then
\begin{equation}\label{eq:sbar2}
(\overline{S}+b-a)^{\varphi_p}=(b-a)+X^{r_1}+X^{-r_1}+X^{r_2}+X^{-r_2}	
\end{equation}	
Since, in relation~\eqref{eq:sbar1}, \(s_0=b-a\), exactly two of the numbers
\(s_1,\ldots,s_{p-1}\) are equal to \(2\), and the remaining
\(p-3\) are equal to \(0\). Moreover, in relation~\eqref{eq:sbar2},
\(s_0=b-a\), exactly four of the numbers
\(s_1,\ldots,s_{p-1}\) are equal to \(1\), and the remaining
\(p-5\) are equal to \(0\). Hence, in both~\eqref{eq:sbar1} and~\eqref{eq:sbar2},
there exist distinct indices \(u,v\in\{1,\ldots,p-1\}\) such that
\(s_u\neq s_v\). Therefore, condition~\eqref{eq:susv2} holds whenever \(p>5\). By assumption \((v)\), \(b-a\neq -4\). It follows that in both~\eqref{eq:sbar1} and~\eqref{eq:sbar2}, \(\sum_{i=0}^{p-1}s_i\neq 0\). So, $
\left|x^h\langle x^p\rangle\cap C\right|
=
\frac{bn}{p(4+b-a)}$ for $0\le h\le p-1$. This completes the proof.\hfill$\blacksquare$



	\end{proof}	

\end{document}
